\documentclass[11pt]{article}

\usepackage[utf8]{inputenc}
\usepackage[T1]{fontenc}
\usepackage[french]{babel}
\usepackage{amsmath,amssymb}
\usepackage[margin=2.3cm]{geometry}
\usepackage{enumitem}
\usepackage{mathtools}

\setlist{itemsep=2pt, topsep=4pt}
\newcommand{\code}[1]{\texttt{#1}}
\newcommand{\fonction}[1]{\code{#1}}

\title{\textbf{TP : entraîner un modèle par descente de gradient}}
\author{Journée des Gymnasiens}
\date{}

\begin{document}
\maketitle

\section*{Objectif}

Comme nous l'avons vu dans l'introduction ce matin, l'apprentissage
automatique consiste à chercher des \emph{paramètres} (des nombres) qui
rendent un \emph{modèle} (une fonction) le plus proche possible des
données. Cette recherche repose sur trois ingrédients que nous avons
déjà rencontrés :

\begin{itemize}
  \item un \textbf{modèle} : une fonction possédant des paramètres
        $p_0, p_1, \dots$, que l'on ajuste pour qu'elle colle aux données ;
  \item des \textbf{données} : des exemples $(x_i, y_i)$ que le modèle doit
        reproduire ;
  \item une fonction de \textbf{coût} (\emph{loss}) : un score qui mesure
        l'écart entre les prédictions du modèle et les cibles $y_i$.
        Il doit être grand quand le modèle se trompe beaucoup et petit quand
        il colle bien aux données.
\end{itemize}

Et nous avons vu enfin comment trouver les paramètres automatiquement :
la \textbf{descente de gradient}, qui à chaque étape modifie les
paramètres dans le sens opposé au gradient du coût :
\[
  p \;\leftarrow\; p - \eta\, \nabla L(p)
\]
où $\eta$ est le \emph{taux d'apprentissage}.

\medskip
Dans ce TP, vous allez vivre cette expérience de l'intérieur, en
programmant vous-même les trois briques. Le fichier \code{src/lab.py}
ne contient que les \textbf{trois fonctions à compléter} :
\fonction{model}, \fonction{loss} et \fonction{grad}. Pour l'instant elles
ne font rien : elles ne renvoient que le symbole \code{...}. Tout le
reste (données, vérifications automatiques, animation d'entraînement) est
déjà écrit dans le lanceur \code{src/lab\_run.py}, qui importe vos trois
fonctions.

\medskip
Après chaque modification, relancez le programme dans VS Code : appuyez
sur \code{F5}. Le programme vérifie d'abord votre
\fonction{model}, puis votre \fonction{loss}, puis votre \fonction{grad}, il lance ensuite l'animation
d'entraînement : elle ne marchera qu'une fois les trois fonctions prêtes.
En cas d'erreur, la vérification affiche la valeur attendue et la valeur
obtenue, ainsi que la taille des tableaux.

\section*{Étape 1 : le modèle \hfill \fonction{model}}

Notre modèle est un polynôme de degré 3, paramétré par une liste
\code{parameters} $= [p_0, p_1, p_2, p_3]$ :
\[
  f(x) \;=\; p_0 \;+\; p_1\,x \;+\; p_2\,x^2 \;+\; p_3\,x^3
  \;=\; \sum_{k=0}^{3} p_k\,x^k .
\]
Le coefficient à la position $k$ multiplie donc $x$ à la puissance $k$.

\begin{enumerate}[label=\textbf{Q\arabic*.}]
  \item Sans calcul : si \code{xs} contient $N$ valeurs, quelle est la
        taille du tableau renvoyé par \fonction{model(parameters, xs)} ?
  \item Implémentez \fonction{model}. L'opérateur \code{**} calcule les
        puissances (\code{xs**2}, \code{xs**3}\ldots), et une boucle
        \code{for k, p in enumerate(parameters)} itère sur chaque
        coefficient $p_k$ avec son rang $k$ : Il faut donc multiplier \code{p} par
        \code{xs**k}, et ensuite sommer les résultats.
  \item Relancez le programme avec \code{F5} : si votre implémentation
        est correcte, le message « \code{Le modèle a l'air bon.} »
        s'affichera. La vérification teste plusieurs valeurs, y compris
        un polynôme de degré plus élevé et plusieurs valeurs de $x$ à la
        fois : votre implémentation doit aussi rester valable quand le
        degré du polynôme ou le nombre de points est différent.
\end{enumerate}

\section*{Étape 2 : le coût \hfill \fonction{loss}}

Comme vu ce matin, nous notons $L$ le \emph{coût} : le score du modèle
pour un jeu de paramètres. Ici (comme dans la plupart des modèles à
sortie numérique), nous prenons la moyenne des carrés des erreurs :
\[
  L(p) \;=\; \frac{1}{N}\,\Bigl( \bigl(f_p(x_1) - y_1\bigr)^2 + \cdots
  + \bigl(f_p(x_N) - y_N\bigr)^2 \Bigr)
  \;=\; \frac{1}{N} \sum_{i=1}^{N} \bigl( f_p(x_i) - y_i \bigr)^2 .
\]

\begin{enumerate}[label=\textbf{Q\arabic*.}, resume]
  \item Que vaut $L$ si le modèle colle parfaitement aux données ? Et si
        le modèle ne donne pas la bonne valeur au moindre point des
        données, que peut-on dire de $L$ ?
  \item Implémentez \fonction{loss(parameters, xs, ys)} : la moyenne des
        carrés des différences entre les prédictions de votre
        \fonction{model} et les cibles \code{ys}. Comme elle fait appel à
        votre \fonction{model}, elle ne peut fonctionner qu'une fois l'étape 1
        terminée. La boucle \code{for x, y in zip(xs, ys):} itère sur les
        paires (valeur de $x$, cible $y$) du jeu de données ; il faut
        ensuite sommer les carrés des différences, puis diviser par la
        taille du jeu de données.
  \item Relancez le programme avec \code{F5} : si vous ne vous êtes pas
        trompé·e, le message « \code{Le coût a l'air bon.} » s'affichera.
\end{enumerate}

\section*{Étape 3 : le gradient \hfill \fonction{grad}}

Nous avions vu ce matin que la descente de gradient demande de connaître
la dérivée partielle $\dfrac{\partial L}{\partial p_i}$ pour chaque
paramètre $p_i$ : de combien $L$ varie si l'on bouge $p_i$ d'une toute
petite quantité. C'est sans doute la partie la plus mathématique du TP,
mais nous allons la calculer \emph{élément par élément} : la collection de
ces dérivées est le \textbf{gradient} $\nabla L$, il a une entrée par
paramètre, et son $i$-ème élément est la dérivée de $L$ par rapport à
$p_i$.

\begin{enumerate}[label=\textbf{Q\arabic*.}, resume]
  \item \emph{Sur papier : les résultats sont donnés, à vous de comprendre
        pourquoi.} Pour un point de données $(x, y)$, posons
        $u = f_p(x) - y$ (l'erreur du point). La règle de dérivation
        des composées donne $\bigl(u^2\bigr)' = 2\,u\,u'$ :
        \begin{enumerate}
          \item Montrez que
                \[
                  \frac{\partial \bigl(u^2\bigr)}{\partial p_i}
                  \;=\; 2\,u \times
                  \frac{\partial f_p(x)}{\partial p_i}\,:
                \]
                pourquoi la cible $y$ disparaît de la dérivée ?
          \item Prouvez que
                \[
                  \frac{\partial f_p(x)}{\partial p_i} \;=\; x^i\,:
                \]
                (dans $f_p(x) = p_0 + p_1\,x + p_2\,x^2 + \cdots$, seul
                le terme en $p_i$ dépend de $p_i$.)
          \item En déduire, en moyennant sur les $N$ points de données :
                \[
                  \frac{\partial L}{\partial p_i} \;=\; \frac{2}{N}
                  \sum_{j=1}^{N} \bigl(f_p(x_j) - y_j\bigr)\,x_j^i\,.
                \]
        \end{enumerate}
  \item Concrètement dans le code : le $i$-ème élément du gradient, à
        savoir $\dfrac{\partial L}{\partial p_i}$, vaut la moyenne sur
        les données de $2\,\bigl(f_p(x_j) - y_j\bigr)\,x_j^i$.
        Implémentez \fonction{grad(parameters, xs, ys)} : la boucle
        \code{for k in range(len(parameters)):} itère les rangs
        $i$, et comme tout à l'heure
        \code{for x, y in zip(xs, ys):} itère sur les paires (valeur de
        $x$, cible $y$) du jeu de données. Pour chaque rang $k$, il faut
        ensuite sommer les éléments, puis diviser par la taille du jeu
        de données.
\end{enumerate}

\section*{Et maintenant : entraînez !}

Maintenant que \fonction{model}, \fonction{loss} et \fonction{grad}
existent, le programme \emph{s'entraîne} tout seul : partant des
paramètres $p = (0, 0, 0, 0)$, il effectue 250 itérations de
$p \leftarrow p - \eta\, \nabla L(p)$ et affiche l'animation
d'entraînement (à gauche les données et les prédictions du modèle, à
droite la courbe du coût).

Faites tourner le programme avec \code{F5} et observez : le polynôme
approxime les points de mieux en mieux. \textbf{Vous avez entraîné
votre premier modèle ! Félicitation !}

\end{document}
